\section[Supplied recursive module responses]{Exact response composition on supplied recursive modules}
\label{sec:op3-recursive-modules}

\paragraph{Status and access distinction.}
The results in this section are \statusproved{} drafts awaiting independent
review. Unlike \cref{thm:op3-canonical-multipartite}, this construction is
given a complete reduced union/join decomposition of the core, all original
core degrees and all private-leaf counts. It does not discover that
decomposition locally. A small response curve and a near-linear
\emph{all-core-edge} construction do not supply an OP3 local-work theorem.
The complete composition and reconstruction are implemented; dense original
obstacle solves are independent validators only.

Use a finite simple connected unweighted graph whose core $C$, of size
$N$, is constructed recursively from singletons by disjoint union and
complete join. A join adds every edge between different children; a union
adds none. A supplied reduced decomposition has at least two children per
internal node and alternating union/join labels along internal edges.
Let $m_C$ be the number of core edges. Each core vertex $i$ has $t_i$
private leaves, and $d_i$ is its original full degree. Initially the
physical seed $v$ lies in $C$; this section does not claim an implemented
arbitrary-seed extension. Keep $0<\alpha<1$, $\lambda>0$, physical
$M=D-\gamma A$, and load $b=e_v-\lambda d$.

For a module $H$, set exterior coordinates to zero and add a uniform
physical field $h\in\mathbb R$ to its core equations, retaining its private
leaves and all original diagonal degrees. Define $F_H(h)$ as the total
core value in this restricted obstacle problem. The field is applied to
the core coordinates only, not to their private leaves. It is an algebraic
parameter; negative fields do not claim to be physical PageRank sources.

\begin{lemma}[Uniform-field module response; \statusproved]
\label{lem:op3-module-response}
Every $F_H$ is continuous, nonnegative, nondecreasing and piecewise affine,
with at most $2|H|$ events. It vanishes for sufficiently negative field.
At a union node $H$ with children $H_j$,
$F_H(h)=\sum_jF_{H_j}(h)$. At a join node, define $G_j$ by
\begin{equation}
 z=t+\gamma F_{H_j}(t),\qquad G_j(z)=F_{H_j}(t),\qquad
 G(z)=\sum_jG_j(z).
 \label{eq:op3-join-input-transform}
\end{equation}
Then $h=z-\gamma G(z)$ is strictly increasing and onto, and
$F_H(h)=G(z)$. On affine pieces the exact transforms are
\begin{align}
 F_{H_j}(t)=at+b
 &\quad\Longrightarrow\quad
 G_j(z)=\frac{az+b}{1+\gamma a},
 &z_*&=t_*+\gamma F_{H_j}(t_*),
 \label{eq:op3-module-child-piece}\\
 G(z)=Az+B
 &\quad\Longrightarrow\quad
 F_H(h)=\frac{Ah+B}{1-\gamma A},
 &h_*&=z_*-\gamma G(z_*).
 \label{eq:op3-module-parent-piece}
\end{align}
Both event transformations preserve order, including on negative intervals.
\end{lemma}
\begin{proof}
The restricted matrix retains original diagonal degrees and is a positive
definite $M$-matrix. Increasing the uniform field increases the load
coordinatewise, so obstacle comparison gives nondecreasing coordinates.
Each core coordinate is born at most once, as is its ordinary private-leaf
group, whose value is $(\gamma u_i-\lambda)_+$. Between these at most
$2|H|$ births the original equations are a fixed nonsingular linear system
with an affine load. This gives the stated affine structure and continuity.
Every core load is nonpositive at sufficiently negative field, and private
leaves have negative load, giving the initial zero piece. A singleton has
exactly the coordinate kernel of \cref{lem:op3-multipartite-scalar}.

The union identity follows from independent children. At a join put
$S=F_H(h)$, $S_j=\sum_{i\in H_j}u_i$, and $z=h+\gamma S$.
Child $j$ sees field $h+\gamma(S-S_j)=z-\gamma S_j$, which gives
\cref{eq:op3-join-input-transform}. Each child inverse is strictly
increasing and onto because $F_{H_j}$ is nondecreasing with finite
piecewise slopes.

It remains to justify the parent inverse without assuming a contraction.
Write $n=|H|$ and $n_j=|H_j|$. On any fixed child face, eliminate its
positive private leaves. The derivative of the remaining positive core
values with respect to its uniform field solves $K w=\mathbf1$.
Here $K$ is a positive definite $M$-matrix whose row sums are at least
$n-n_j$. Indeed, every cross-child core edge contributes its original
diagonal unit and no child off-diagonal entry. Internal positive core
neighbors leave a nonnegative $(1-\gamma)$ contribution; internal zero
core neighbors leave a full diagonal unit. Private leaves leave
$t_i-\gamma^2t_i^{+}\geq0$, and original core neighbors outside $H$
only improve the margin. Therefore
$K^{-1}\mathbf1\leq\mathbf1/(n-n_j)$ and
\[
 0\leq F'_{H_j}\leq\frac{n_j}{n-n_j},\qquad
 \gamma G'_j\leq
 \frac{\gamma n_j}{n-n_j+\gamma n_j}<\frac{n_j}{n}.
\]
This argument applies to all affine intervals, including negative ones.
Summing gives $\gamma G'<1$ with a fixed positive margin for this finite
node. Thus $z-\gamma G(z)$ is strictly increasing and onto. The affine
inverse identities and event positions follow by substitution. They also
show directly that a parent creates no more events than its children
together. Simultaneous events are merged at their common continuous value.
The recovered child equations are precisely the original restricted
obstacle conditions, so uniqueness identifies the composed response.
\end{proof}

\begin{lemma}[All intermediate module sizes are paid by distinct edges; \statusproved]
\label{lem:op3-module-edge-charge}
In a supplied reduced decomposition,
\begin{equation}
 \sum_{H\text{ internal}}|H|\leq N+4m_C.
 \label{eq:op3-module-edge-charge}
\end{equation}
No balance or logarithmic-height assumption is needed.
\end{lemma}
\begin{proof}
For a join node of size $n\geq2$ with at least two children, its cross-child
edge set has size $\sum_{j<k}n_jn_k\geq n-1\geq n/2$.
Every core edge belongs to exactly one such set, at the lowest common
ancestor of its endpoints. Hence the sum of join-node sizes is at most
$2m_C$. Each non-root union node has a join parent. The total sizes of all
union children of a particular join are at most that parent's size.
Consequently the sum of non-root union-node sizes is at most the sum of
join-node sizes. A union root, if present in a restricted core instance,
adds at most $N$. This proves \cref{eq:op3-module-edge-charge}.
An alternating comb can have quadratically many vertex records, but it
can also have quadratically many core edges; the former count alone is not
an excessive-work witness in this supplied-input model.
\end{proof}

\begin{theorem}[Supplied recursive response construction; \statusproved]
\label{thm:op3-supplied-modules}
Given a reduced postorder decomposition, all original core degrees and
private-leaf counts, the implemented bottom-up construction stores every
module response and the required child transforms in
\begin{equation}
 O((N+m_C)\log(2+N))\ \text{exact-real word work},\qquad
 O(N+m_C)\ \text{words of space}.
 \label{eq:op3-supplied-module-work}
\end{equation}
For any supplied field $h$, it recovers all core values in
$O(N\log(2+N))$ further work. Each private-leaf group's common value then
costs constant work per parent. Explicit original leaf output, if requested,
additionally charges enumeration of the relevant original incidences and
each emitted coordinate. These are all-core-input bounds, not
$\operatorname{cvol}(U)$ bounds.
\end{theorem}
\begin{proof}
Read the supplied tree and kernel metadata once, checking its postorder,
unique-parent, alternating-label and vertex-cover properties with charged
dictionary operations. This does not check that an unread ambient graph
has the supplied decomposition; that correspondence is an input promise.
The reduced tree has $O(N)$ nodes and child references.

Build each singleton kernel once. At a union, merge the children's sorted
event streams with one heap and one current affine pair per child. At a
join, first transform each child event using
\cref{eq:op3-module-child-piece}; merge those streams and transform each
combined event using \cref{eq:op3-module-parent-piece}. Combine equal
thresholds before advancing. Every input event has one heap insertion,
one removal and constant many coefficient operations. Temporary current
pieces, event records and heap entries are charged. Retain child curves
and join-input transforms for recovery, rather than treating their copies
as free.

The response lemma bounds the sum of input stream lengths at a node by
twice its size. Thus all transforms, copies, heap work, retained curves
and tree validation cost the claimed amount by
\cref{lem:op3-module-edge-charge}. Peak space includes all retained
histories, the current merge workspace and dictionaries. They obey the
same linear all-core-edge record bound.

To recover, query the root curve by binary search and traverse the tree
once with a stack. A union passes its field to its children. A join uses
its known mass $S$ to form $z=h+\gamma S$, queries each stored $G_j(z)$,
and passes $z-\gamma G_j(z)$ and that child mass to the child. Each node
is visited once and each child mass requires one binary search. A
singleton returns its coordinate response, and private leaves follow from
their scalar formula. No full vector is written between parent
construction steps. The algorithm returns every core value, including
zeros; its input charge is therefore not support-local.
\end{proof}

\paragraph{Measured and open.}
The full \nolinkurl{recursive_module_response.py} audit has 1,716 executions
on all 53 eligible connected graph-atlas cores through six vertices, every
core seed, two unequal-leaf patterns and three parameter pairs. It checks
15,102 module curves through 50,250 complete original affine-interval
certificates and 50,250 original event or zero-field comparisons, including
20,063 negative-field interval samples. Each interval certificate derives
the original affine coordinate vector independently, verifies equality
coefficients and inequality signs over the entire interval including
unbounded endpoints, and compares its core-mass coefficients to the stored
curve. Ten additional balanced and alternating-comb instances through 128
core vertices check reconstruction and the edge ledger.

The supplied full-graph decomposition recognizer and dense validators are
audit-only. The source of a local module representation remains
\statusopen{}. In particular, the distinct-edge charge uses edges between
inactive core vertices too, and a positive pair inside one deep module need
not enumerate the whole core. One cannot import
\cref{eq:op3-supplied-module-work} into the local multipartite theorem by
declaring a decomposition oracle free. General OP3 remains \statusopen{}.

\subsection{Persistent composition without an all-core-edge factor}
\label{sec:op3-persistent-modules}

The explicit merge above has a stronger implementation using the immutable
affine curve structure already developed in
\cref{lem:op3-persistent-primitives}. The supplied decomposition remains an
essential assumption. The improvement concerns algebra on that compressed
input, not the cost of constructing it from local graph queries.

\begin{lemma}[Convex zero-left-ray curves and legal shears; \statusproved]
\label{lem:op3-module-convex-shears}
Every module curve $F_H$ and transformed child curve $G_j$ is convex and
has a zero left ray. They have representations
\begin{equation}
 F(h)=\sum_{\ell=1}^{k}c_\ell(h-\tau_\ell)_+,
 \qquad c_\ell\geq0,
 \label{eq:op3-module-hinges}
\end{equation}
with $k\leq2|H|$. The child and parent join transforms are the respective
whole-curve shears
\begin{equation}
 (x,y)\longmapsto(x+\gamma y,y),\qquad
 (x,y)\longmapsto(x-\gamma y,y).
 \label{eq:op3-module-shears}
\end{equation}
Each preserves strict horizontal point order on its applicable curve.
Their persistent implementations take constant work and new records.
Adding a zero-left-ray curve with $k$ knots to another takes
$O(k\log(2+N))$ work and new records without traversing the other curve.
\end{lemma}
\begin{proof}
Singleton kernels are convex and start at zero. Unions preserve convexity.
At a child transform the new slope is $a/(1+\gamma a)$, an increasing
function of $a\geq0$. At the parent transform it is $A/(1-\gamma A)$,
an increasing function of $A$ on the domain
$\gamma A<1$ proved in \cref{lem:op3-module-response}. The event
transformations preserve order, so this inductively proves convexity and
the hinge representation. Both shears fix the zero ray; their horizontal
derivatives are positive throughout the applicable slope range.

Store the curve knots in the same immutable affine-tagged AVL tree as in
\cref{lem:op3-persistent-primitives}, with both ray slopes in its header.
The present left slope is zero. This does not invalidate horizontal
search, interpolation, insertion, suffix shears or affine tagging; no
vertical inverse query is used. The search tree's stored heights and
relative horizontal order stay valid under \cref{eq:op3-module-shears}.
The same constant-size root copies and path-copying proof therefore apply.
Only horizontal monotonicity and a nonzero affine determinant are needed
for these operations; the stricter inverse-query condition is not imported.

Stream the smaller curve's knots in order, obtaining the nonnegative
slope jumps in \cref{eq:op3-module-hinges}. For each jump insert its knot
in the recipient if needed, then apply the continuous suffix shear
$(x,y)\mapsto(x,y+c_\ell(x-\tau_\ell))$ above that threshold. Each step
copies one search path and tags whole off-path subtrees. It changes only
$O(\log(2+N))$ nodes and leaves every prior version unchanged. A streaming
iterator visits each smaller-curve knot once and retains only its search
stack. The zero initial ray requires no initial linear-term addition.
\end{proof}

\begin{theorem}[Persistent solve on a supplied recursive core; \statusproved]
\label{thm:op3-persistent-modules}
Under the supplied-input assumptions of \cref{thm:op3-supplied-modules},
all module curves, required transformed child versions, and one final
recovery of all core coordinates can be constructed in
\begin{equation}
 O(N\log^2(2+N))
 \label{eq:op3-persistent-module-work}
\end{equation}
exact-real word work and allocated words. There is no $m_C$ factor and no
polynomial arithmetic-operation dependence on $1/\alpha$ or $1/\lambda$.
The output and locality qualifications of
\cref{thm:op3-supplied-modules} still apply.
\end{theorem}
\begin{proof}
Compute physical module sizes from the supplied postorder tree. At a
union, retain the curve of a largest-size child as the initial sum and
merge each other child by \cref{lem:op3-module-convex-shears}. At a
join, first apply the positive shear to every child using one persistent
root copy per child, retain a largest-size child's transformed curve,
merge the other transformed curves, then apply the negative shear to
their sum. Largest means physical core count, not the number of distinct
knots, which can shrink under ties. Store every module root and each
transformed child root for reconstruction.

If a child is not the chosen largest child, its size is at most half
the parent's size. Each physical vertex can therefore belong to such a
child at most $\lfloor\log_2N\rfloor$ times on its ancestor path.
The sum of all merged non-largest child sizes is at most
$N\lfloor\log_2N\rfloor$. Each such curve has at most twice its physical
size in knots. This pays for every streamed knot, inserted breakpoint,
suffix update, lazy push, AVL rotation and immutable copied record in
\cref{eq:op3-persistent-module-work}. There are only $O(N)$ tree nodes,
child references, retained version headers, whole-curve maps and singleton
events. Their metadata reads, comparisons and allocations fit the bound.
The bound counts all allocated nodes, so it includes old recovery
versions and even intermediate records; it does not require freeing
history to make space appear small.

The downward traversal is unchanged, querying stored transformed children
at each join and original child curves at each union. Horizontal searches
carry composed affine maps without modifying or copying tree nodes. There
are $O(N)$ queries and output core coordinates, each query costing
$O(\log(2+N))$ work. Conditional original obstacle equations from
\cref{lem:op3-module-response} prove the recovered response exactly.
Private-leaf group values are obtained once per parent; explicit original
leaf enumeration remains an additional charged output operation.

The input can represent $\Theta(N^2)$ core edges implicitly. This is why
an $m_C$-independent bound is consistent with the access model. Reading an
explicit graph or finding its decomposition is not silently included as
free preprocessing. No local support-volume theorem follows from the
compressed-input theorem alone.
\end{proof}

\paragraph{Implementation and independent verification.}
\nolinkurl{persistent_module_response.py} imports the existing immutable
AVL implementation without changing it. Its new sum routine handles
zero left rays and streams only non-largest child curves. The audit checks
every retained module and transformed child version after all parents have
been built. It compares complete affine pieces against the explicit module
engine's independently certified original-matrix intervals, checks AVL
height/size invariants, and recovers coordinates at events, interval
interiors, zero and both unbounded rays. Validation uses a separate observer
so reads of old versions do not masquerade as construction work. One final
production recovery and all actual construction allocations are counted
separately. The saved full audit records its exact scope and code hashes.
It passes 1,716 original graph cases, 50,268 retained affine-piece matches,
22,530 saved curve versions and 100,500 module-coordinate recoveries,
alongside the 50,250 independent original interval certificates. Eighteen
additional comb, clique and star diagnostics reach 512 core vertices.
The 512-vertex alternating comb has 65,536 core edges and allocates
38,598 immutable AVL nodes; this is a measured construction count,
not a local-discovery benchmark.
